\usepackage[spanish,es-tabla]{babel}
\decimalpoint

\usepackage[a4paper]{geometry} % Tamaño del papel
\usepackage{graphicx}               % Para incluir imágenes
\usepackage{subcaption}
\usepackage{amsmath}                % Para el manejo de matemáticas
\usepackage{amsthm}
\usepackage{multirow}					% Combinar columnas
\usepackage{url}					% Dar formato a urls				
\usepackage{array}					% Para ajustar el texto en la celda
\usepackage{xcolor}
\usepackage{float}
\usepackage{comment} 
\usepackage{amssymb}		% Simbolos matematicos
\usepackage{fancyhdr}  % Para encabezados personalizados
\newtheorem{definition}{Definición}
\setlength{\parskip}{10pt}

% Configurar encabezados
\pagestyle{fancy}
\fancyhf{}
\fancyhead[LE,RO]{\small\leftmark}    % Muestra el capítulo en lugar de la sección
\fancyhead[RE,LO]{\small\thepage}      % Número de página
\setlength{\headheight}{13.6pt}

\setcounter{secnumdepth}{3}